\section{Methodology}
\label{sec:failure_forecasting_methodology}

    \subsection{Data}

\textit{Bicing} has made available for this use case two bike-sharing data sets: a trips dataset and a maintenance dataset (Appendix~\ref{app_paper3:section_data_exploration}).

\begin{itemize}
    \item The trips dataset is described in detail in Chapter~\ref{ch:bicing_mobility_patterns}. For the mobility analysis presented in that section, only trips data from 2022 was used to minimize the impact of the COVID-19 pandemic and to reflect the increasing use of \acp{e-b}. In contrast, the maintenance analysis and the \ac{pm} modeling utilized the complete trips dataset available (from April 2019 to December 2022).

    \item Maintenance data is comprised of a total of 310,000 \acp{mo}, which correspond to various bicycle repairs types executed between September 1st, 2020, and January 1st, 2023 (i.e., 2 years and 4 months). Each \ac{mo} record provides the following information: \ac{mo} identifier, date, category, subcategory, bike identifier, and bike model (\ac{m-b} or \ac{e-b}).
    There exist a total of 12 categories and 87 subcategories, encompassing a wide array of actions such as cleaning, greasing, adjusting, or changing bike parts. Categories range from fewer than 100 interventions to 120,000, with the majority of them falling into the brake and wheel categories (66\%). Likewise, subcategories also reveal a wide range of counts, with only 10\% of the repair typologies surpassing 10,000 \acp{mo}.
  \end{itemize}

    \subsection[MOs processing]{\acp{mo} processing} 
    \label{section_mos_processing}

First, the target bike parts for this study needed to be selected. To ensure complete objectivity, subjective \ac{mo} types such as cleaning or greasing were excluded, and the focus was placed on the replacement of bike parts. Also, failures related to wheel tubes were excluded due to their high level of randomness.

\ac{mo} data records specific dates for bike part repairs, however, this data structure is not optimal for conducting \ac{sa}, which typically models the time to an event. To facilitate the use of \ac{sa}, \acp{mo} were transformed into \ac{mo} units. These units represent the time elapsed between two consecutive repairs for the same bike and bike component, and thus, it is the time period in which the bike part under study is operational. Consequently, when working with \ac{mo} units, bike part survival information at the beginning and end of the data set could be lost for each bike. To prevent this information loss, the time from the start of the data set to the first \ac{mo} is considered as one \ac{mo} unit, and the time from the last repair of the bike to the end of the data set is regarded as another \ac{mo} unit. It's worth noting that even though \ac{mo} units may originate from the same bike, each one has been treated as an independent entity.

One key challenge in predicting the occurrence of an event is the presence of censored data, which refers to incomplete information about survival times~\cite{Leung1997, Gijbels2010}. Distinct types of censoring exist, including right-censoring (the event has not occurred by the end of the study), left-censoring (the event occurred before the study started), and interval-censoring (the event occurred within a specific time interval). Bicing maintenance data contains left- and right-censored data, since the first and last \ac{mo} units of each bike and repair typology are incomplete. In the first unit, it is not possible to know when this bike part started working, and in the last, when this bike part finally will break. Since \ac{sa} can effectively utilize uncensored and right-censored data to estimate the survival curves and generate predictions, only the left-censored \ac{mo} units were discarded in the training and prediction phases. Specifically, this involves excluding the initial \ac{mo} unit of each bike for the target bike part. Additionally, units without trips and the ones in which \ac{m-b} were upgraded to \ac{e-b} were also discarded.

The data sets used for the survival models follow a format where each row encapsulates all the pertinent information about one subject. Within each entry, there are the details about the subject's survival duration, event occurrence, and aggregated covariate values crucial for the survival function estimation. \ac{mo} units, with their defined start and end dates, facilitate the calculation of covariates that describe the utilization patterns of a bike part and its surrounding environmental factors. In this way, various covariates domains have been incorporated to the model inputs:

\begin{itemize}
\item \textbf{Weather}: weather data was analyzed to incorporate the surrounding environmental conditions of the \ac{mo} unit. This involved considering the daily average temperature (in ºC), total daily precipitation (in mm), average wind direction (in degrees), mean wind speed (in km/h), and average atmospheric pressure (in hPa). After exploring this data and its connection to bike part failures, it was determined that the primary influential factors were the daily average temperature and the mean atmospheric pressure.
\item \textbf{Bike usage}: various metrics were computed, including daily distance traveled (in meters), daily positive and negative inclinations (in meters), and the mean daily speed (in km/h). These calculations were based on the previously collected trip routes and inclinations between each couple of stations. Next, their cumulative versions were examined, and strong correlations were identified. As a result, only two variables were kept: the cumulative daily distance and the mean speed of the \ac{mo} unit.
\item \textbf{Bike model}: due to potential variations in survival curves across the two bicycle models, a binary variable was introduced to distinguish between electric (1) and mechanical (0) bike models.
\item \textbf{Count of repairs for other bike parts during the target \ac{mo} unit}: when analyzing a specific bike part, it becomes pertinent to evaluate the frequency of replacements for another bike part. This approach allows to gain valuable insights; for instance, understanding the number of wheel tubes replaced could aid in predicting potential replacements for components such as tires or wheel rims. Following an examination of the correlation and the variation inflation factor of these repair counts, the following \ac{mo} subcategories were selected: brake tension adjustment, the replacement of the front and rear wheel tubes, and the replacement of the front wheel cover.
\end{itemize}

    \subsection{Models}
To predict the survival time of the bike components, the following statistical and machine learning models have been employed:

\begin{enumerate}
\item \textbf{\acl{cph}} (\acs{cph})\acused{cph} model~\cite{Cox1972} (lifelines' implementation~\cite{Lifelines2021}): is a semi-parametric approach that enables to evaluate how covariates influence the hazard rate of an event as time progresses. This model relies on an underlying assumption known as the proportional hazard assumption, which asserts that the relative risk between two distinct groups remains constant over time. Meeting this assumption simplifies the analytical process and enhances the interpretability of outcomes. Although evaluating the assumption holds theoretical importance for attaining a meaningful interpretation of covariates, adherence to this assumption might not be imperative.~\cite{Stensrud2020} noted that when dealing with a sufficiently large sample size, even minor deviations from the assumption may show up. Furthermore, when the main goal is survival prediction, there is no need to test the proportional hazard assumption, since the main objective is to maximize an score~\cite{lifelines_ph_assumption}. Consequently, for the purpose of this study, the analysis will prioritize predictive accuracy over strict adherence to the assumption.

\item \textbf{\acl{mtlr}} (\acs{mtlr})\acused{mtlr} model~\cite{Yu2021} (pysurvival's implementation~\cite{pysurvival}): stands as an alternative to \ac{cph} when the assumption of proportional hazards does not hold. \ac{mtlr} model relies on a sequence of logistic regression models constructed across distinct time intervals. This allows the estimation of the probability associated with the occurrence of the event of interest within each interval. Consequently, the initial step involves specifying the desired number of intervals, with the present use case opting for the number of days in the maintenance data.

\item \textbf{\acl{csf}} (\acs{csf})\acused{csf} model~\cite{Wright2017} (pysurvival's implementation~\cite{pysurvival}): is a machine learning model that extends Random Forest ensembles to effectively handle right-censored data. Therefore, this survival model allows to appropriately model data with nonlinear relationships and censoring.

\item \textbf{\acl{deepsurv}} (\acs{deepsurv})\acused{deepsurv} model~\cite{Katzman2018} (pysurvival's implementation~\cite{pysurvival}): is an improved version of the \ac{cph} model that brings in elements of deep learning to its core structure. This enhancement enables the model to better capture complex patterns while still being able to handle censored data effectively.
\end{enumerate}

    \subsection{Hyperparameters optimization}
The previous models possess distinct architectures. Consequently, the hyperparameter optimization process has varied based on the specific model.

For the purpose of identifying the optimal hyperparameters, the \ac{sa} dataset of each bike component was partitioned into three subsets: training (60\%), validation (20\%), and test (20\%). Multiple hyperparameter configurations were trained on the training set and assessed with the validation set and the RMSE metric to identify the most favorable combination. Once determined, the model was trained using both the training and validation sets, and predictions based on the test set were generated to asses the final accuracy of the model. Furthermore, although model validation was primarily assessed using the RMSE metric, all data subsets were also evaluated for RMSE, R\textsuperscript{2} and MAPE to assess accuracy comprehensively (Appendix~\ref{app_paper3:section_accuracy_measures}). RMSE and MAPE interpretation is very straightforward; the lower, the better. However R\textsuperscript{2} works in the contrary direction since higher values, correspond to better predictions.

For the \ac{cph} model, the sole hyperparameter subjected to optimization was the baseline estimation method (breslow, spline, or piecewise), while the remaining parameters were set to their default values. In the case of \ac{mtlr}, \ac{csf}, and \ac{deepsurv}, the hyperparameter search process was automated using the Optuna framework~\cite{Optuna2019}. A total of 200 trials were conducted for each model, employing the TPEsampler class for hyperparameter sampling and the MedianPruner class to halt unpromising combinations. For \ac{mtlr}, the optimized hyperparameters encompassed the learning rate (ranging between 1e-5 and 1e-3), initialization method (orthogonal or glotorot\_uniform), and optimizer (adam, adamaz, or sgd). The \ac{csf} optimization encompassed the number of trees (ranging between 10 and 100 in increments of 10), maximum tree depth (between 2 and 10), and minimum node size (ranging between 10 and 50 in steps of 5). In the case of \ac{deepsurv}, the optimization process considered initialization method (orthogonal or glotorot\_uniform), optimizer (sgd or adam), learning rate (ranging between 1e-5 and 1e-2), number of epochs (ranging between 50 and 500), L2 regularization (ranging between 0 and 1e-2), and the inclusion of batch normalization or dropout with a value of 0.5 (True or False for each one).
